% !TEX root = ../main.tex
\corrigesChapitre{12}{Stratégies de synthèse organique}
\corrigesSouples

\begin{blocLivre}{cor-12-1}
\begin{corrige}{1}{Les étapes d'un protocole}
\demarche{Transformation, isolement, purification, identification : chaque phrase du protocole sert l'une des quatre.}
\begin{questions}
\item (1) et (2) : \res{\text{transformation}} ; (3) : \res{\text{isolement}} ; (4) : \res{\text{purification}} ;
(5) : \res{\text{identification}} (contrôle de pureté).
\item L'acide sulfurique est le \res{\text{catalyseur}} ; la pierre ponce régularise l'ébullition et évite les
soubresauts.
\item À reflux, les vapeurs sont condensées et retombent dans le ballon : on peut chauffer, donc accélérer, sans
perdre ni réactifs ni produit volatils.
\item Le produit \ce{CH3-COO-CH2-CH2-CH(CH3)2} est un \res{\text{ester}}.
\end{questions}
\end{corrige}
\end{blocLivre}
\begin{blocLivre}{cor-12-2}
\begin{corrige}{2}{Le montage à reflux}
\demarche{Suivre le trajet des vapeurs : elles montent du ballon, se condensent, retombent.}
\begin{questions}
\item A : \res{\text{ballon}} contenant le mélange réactionnel ; B : \res{\text{réfrigérant à eau}} ; C :
\res{\text{entrée d'eau froide}} ; D : \res{\text{sortie d'eau}} ; E : \res{\text{chauffe-ballon}}.
\item L'eau entre \res{\text{en bas}} et sort en haut : le réfrigérant reste entièrement rempli et l'eau la plus
froide rencontre les vapeurs les plus chaudes.
\item Les vapeurs qui montent se condensent sur la paroi froide et \res{\text{retombent dans le ballon}} : la matière
reste dans le montage, même à ébullition prolongée.
\item Le \res{\text{réfrigérant}} : il empêche les vapeurs, souvent inflammables, de se répandre dans la pièce.
\end{questions}
\end{corrige}
\end{blocLivre}
\begin{blocLivre}{cor-12-3}
\begin{corrige}{3}{Calculer un rendement}
\demarche{La quantité maximale d'ester est celle du réactif limitant ; le rendement compare masses obtenue et maximale.}
\begin{questions}
\item $\res{\ce{CH3COOH + C2H5OH <=> CH3COOC2H5 + H2O}}$.
\item $n_{\max}=n(\text{acide})=\qty{0,100}{mol}$, soit $m_{\max}=0{,}100\times88{,}1=\res{\qty{8,8}{g}}$.
\item $\eta=\dfrac{5{,}9}{8{,}81}=\res{0{,}67}$, soit 67~\%.
\item La transformation est \res{\text{non totale}} ($K=4$, chapitre~7) et une partie du produit est perdue lors de
l'isolement et de la purification.
\end{questions}
\end{corrige}
\end{blocLivre}
\begin{blocLivre}{cor-12-4}
\begin{corrige}{4}{Substitution, addition ou élimination ?}
\demarche{Addition : une liaison multiple s'ouvre ; élimination : elle se crée ; substitution : un groupe en remplace un autre.}
\begin{questions}
\item A : la double liaison s'ouvre et fixe \ce{H} et \ce{OH} : \res{\text{addition}}. B : \res{\text{substitution}}.
C : une double liaison se forme avec départ d'eau : \res{\text{élimination}}. D : \res{\text{addition}} de \ce{H2}.
\item Dans B, l'atome de brome est \res{\text{remplacé}} par le groupe hydroxyle \ce{-OH} ; l'ion bromure part.
\item La polymérisation du chapitre~11 enchaîne des \res{\text{additions}} sur la double liaison C$=$C, comme en A
et D.
\end{questions}
\end{corrige}
\end{blocLivre}
\begin{blocLivre}{cor-12-5}
\begin{corrige}{5}{Chaîne ou groupe ?}
\demarche{Comparer le squelette carboné et le groupe caractéristique avant et après.}
\begin{questions}
\item a. Le squelette \ce{C2} est conservé, l'alcool devient acide : modification \res{\text{de groupe}}.
b. Le squelette s'allonge indéfiniment : modification \res{\text{de chaîne}}.
c. \ce{-OH} remplacé par \ce{-Br}, squelette intact : modification \res{\text{de groupe}}.
d. L'acide et l'alcool deviennent un ester : modification \res{\text{de groupe}}.
\item La transformation \res{\text{c}} est une substitution : un groupe y est remplacé par un autre.
\end{questions}
\end{corrige}
\end{blocLivre}
\begin{blocLivre}{cor-12-6}
\begin{corrige}{6}{La synthèse de l'aspirine}
\demarche{Le réactif limitant est l'acide salicylique ; le rendement compare la masse pesée à la masse maximale.}
\begin{questions}
\item Transformation : reflux 20~min ; isolement : cristallisation sur glace et filtration sur Büchner ;
purification : recristallisation et séchage ; identification : température de fusion.
\item Pour que l'acide salicylique soit \res{\text{limitant}} et soit transformé au maximum : l'excès d'anhydride
déplace l'état final dans le sens direct.
\item $n_{\max}=\dfrac{5{,}00}{138{,}1}=\qty{3,62e-2}{mol}$, d'où
$m_{\max}=\num{3.62e-2}\times180{,}2=\res{\qty{6,52}{g}}$ et $\eta=\dfrac{4{,}2}{6{,}52}=\res{0{,}64}$, soit 64~\%.
\item La température de fusion mesurée (135~°C) est quasi égale à la valeur de référence (136~°C) : le solide est bien
de \res{\text{l'aspirine, pratiquement pure}}.
\end{questions}
\end{corrige}
\end{blocLivre}
\begin{blocLivre}{cor-12-7}
\begin{corrige}{7}{Choisir les bonnes conditions}
\demarche{Vitesse : facteurs cinétiques (chapitre 4) ; rendement : état d'équilibre et critère d'évolution (chapitre 7).}
\begin{questions}
\item a. Catalyseur : augmente \res{\text{la vitesse seulement}}. b. Température : facteur cinétique, augmente
\res{\text{la vitesse}} (l'état final, fixé par $K$, n'est pas amélioré). c. Excès d'éthanol : déplace l'état final,
augmente \res{\text{le rendement}} (et accélère un peu, les chocs étant plus fréquents). d. Éliminer l'eau : augmente
\res{\text{le rendement}}, la transformation devenant quasi totale.
\item Un catalyseur accélère l'aller et le retour sans changer $K$ : \res{\text{l'état final est inchangé}}
(chapitre~7).
\item Retirer l'eau fait \res{\text{diminuer }Q_\mathrm{r}}, qui reste inférieur à $K$ : le système continue d'évoluer
dans le sens direct.
\end{questions}
\end{corrige}
\end{blocLivre}
\begin{blocLivre}{cor-12-8}
\begin{corrige}{8}{L'économie d'atomes}
\demarche{EA = masse molaire du produit visé divisée par la somme des masses molaires des réactifs.}
\begin{questions}
\item $\mathrm{EA}=\dfrac{M(\text{produit visé})}{\sum M(\text{réactifs})}\times100$, nombres stœchiométriques
compris.
\item Voie A : $\dfrac{60{,}1}{42{,}1+18{,}0}=\res{100~\%}$ ; voie B :
$\dfrac{60{,}1}{123{,}0+40{,}0}=\dfrac{60{,}1}{163{,}0}=\res{37~\%}$.
\item Le bromure de sodium \ce{NaBr} : près des deux tiers de la masse engagée finit en \res{\text{déchet salin}} à
traiter.
\item Dans une addition, les réactifs se combinent sans rien rejeter : \res{\text{tous leurs atomes}} se retrouvent
dans le produit.
\end{questions}
\end{corrige}
\end{blocLivre}
\begin{blocLivre}{cor-12-9}
\begin{corrige}{9}{Le rendement d'une synthèse multiétapes}
\demarche{Chaque étape ne conserve qu'une fraction de la précédente : les fractions se multiplient.}
\begin{questions}
\item L'étape 2 part du produit de l'étape 1 : les fractions conservées se multiplient,
$\eta=\eta_1\times\eta_2\times\eta_3$.
\item $\eta=0{,}85\times0{,}90\times0{,}70=\res{0{,}535}$ : il reste
$0{,}50\times0{,}535=\res{\qty{0,268}{mol}}$ de produit final.
\item Améliorer \res{\text{l'étape à 70~\%}}, la plus faible : le gain, facteur $95/70$, est le plus grand. Nouveau
rendement : $0{,}85\times0{,}90\times0{,}95=\res{0{,}73}$.
\item Une seule étape à 70~\% donnerait $\eta=\res{70~\%}$, mieux que les trois étapes réunies (53,5~\%) : réduire le
nombre d'étapes est le levier le plus puissant.
\end{questions}
\end{corrige}
\end{blocLivre}
\begin{blocLivre}{cor-12-10}
\begin{corrige}{10}{Pourquoi protéger ?}
\demarche{Chaque groupe amine libre peut attaquer chaque groupe carboxyle libre.}
\begin{questions}
\item Chaque molécule porte les deux groupes : tout amine peut se coupler à tout carboxyle, y compris entre molécules
identiques. On obtient \res{\text{Gly--Gly, Gly--Ala, Ala--Gly et Ala--Ala}}.
\item Protéger \res{\text{l'amine de Gly}} et \res{\text{le carboxyle d'Ala}} ; coupler : seul le carboxyle libre de
Gly peut réagir avec l'amine libre d'Ala, formant Gly--Ala ; enfin déprotéger.
\item Les groupes bloqués doivent être rendus au dipeptide : sans déprotection, le produit n'aurait ni son amine ni
son carboxyle d'extrémité, et ne pourrait pas être allongé ou utilisé.
\item Deux protections et des déprotections s'ajoutent au couplage : le rendement global, \res{\text{produit des
rendements}} de toutes ces étapes, diminue d'autant.
\end{questions}
\end{corrige}
\end{blocLivre}
\begin{blocLivre}{cor-12-11}
\begin{corrige}{11}{La juste durée du reflux}
\demarche{La montée traduit la cinétique ; le palier traduit l'équilibre chimique.}
\begin{questions}
\item Le rendement croît puis atteint un \res{\text{palier à 78~\%}} : l'état d'équilibre est atteint, la composition
n'évolue plus ($Q_\mathrm{r}=K$).
\item Environ \res{\qty{50}{min}} : le rendement vaut alors 77~\%, à un point du palier ; chauffer davantage consomme
de l'énergie sans presque rien produire de plus.
\item Avec un catalyseur, la montée est plus rapide mais le palier \res{\text{garde la même valeur}} : la courbe verte
rejoint 78~\% plus tôt (voir figure).
\item Déplacer l'état final : \res{\text{excès d'un réactif}} ou \res{\text{élimination de l'eau}} au fur et à mesure.
\end{questions}
\schema[corrige]{12-ex11}
\end{corrige}
\end{blocLivre}
\begin{blocLivre}{cor-12-12}
\begin{corrige}{12}{Deux routes vers l'ibuprofène}
\demarche{EA = M(produit)/somme des M(réactifs) ; les sous-produits représentent le complément.}
\begin{questions}
\item Boots : $\dfrac{206{,}3}{514{,}5}=\res{40{,}1~\%}$ ; BHC : $\dfrac{206{,}3}{266{,}8}=\res{77{,}3~\%}$ :
les valeurs annoncées sont confirmées.
\item Par tonne d'ibuprofène : Boots rejette $\dfrac{514{,}5-206{,}3}{206{,}3}=\res{\qty{1,5}{t}}$ de sous-produits ;
BHC, $\dfrac{266{,}8-206{,}3}{206{,}3}=\res{\qty{0,29}{t}}$ : cinq fois moins.
\item Le rendement global est le \res{\text{produit}} des rendements des étapes : passer de six à trois étapes
multiplie mécaniquement le rendement global, en plus de simplifier le procédé.
\item \res{\text{Catalyse}} avec catalyseurs récupérés, \res{\text{économie d'atomes}} élevée (moins de déchets),
réduction du nombre d'étapes.
\end{questions}
\end{corrige}
\end{blocLivre}
\begin{blocLivre}{cor-12-13}
\begin{corrige}{13}{Assembler un dipeptide sans fausse note}
\demarche{Le couplage amide doit n'avoir qu'un seul choix : bloquer les deux groupes qui ne doivent pas réagir.}
\begin{questions}
\item Le couplage crée un groupe \res{\text{amide}} \ce{-CO-NH-} : le dipeptide relève de la famille des amides.
\item Sans protections, chaque amine libre peut attaquer chaque carboxyle libre : on obtiendrait les
\res{\text{quatre dipeptides}} (et des enchaînements plus longs), au lieu du seul Ala--Gly.
\item Dans Ala--Gly, c'est le \res{\text{carboxyle de l'alanine}} qui forme la liaison : il doit rester libre, et son
amine doit être bloquée ; symétriquement, la glycine doit garder son amine libre et son carboxyle protégé.
\item $\eta=0{,}90^4=\res{0{,}66}$ : à peine deux tiers du réactif aboutit au dipeptide.
\item Avec $n$ couplages, le rendement varie comme $0{,}90^n$ : pour trente étapes, il resterait 4~\% du produit.
Chaque étape doit donc être quasi quantitative, ce qui explique la sophistication des synthèses peptidiques.
\end{questions}
\end{corrige}
\end{blocLivre}
\begin{blocLivre}{cor-12-14}
\begin{corrige}{14}{La synthèse du paracétamol}
\demarche{Le 4-aminophénol est limitant ; la masse maximale s'obtient mole à mole.}
\begin{questions}
\item L'amine est transformée en \res{\text{amide}} : le paracétamol porte un groupe amide et conserve le groupe
hydroxyle du cycle.
\item Transformation : ajout d'anhydride et reflux ; isolement : eau glacée puis filtration sur Büchner ;
purification : recristallisation et séchage ; identification : température de fusion.
\item $n_{\max}=\dfrac{5{,}45}{109{,}1}=\qty{5,00e-2}{mol}$, soit
$m_{\max}=\num{5.00e-2}\times151{,}2=\res{\qty{7,55}{g}}$ ; $\eta=\dfrac{6{,}0}{7{,}55}=\res{0{,}79}$, soit 79~\%.
\item Sans recristallisation, on garderait environ $\dfrac{6{,}0}{0{,}90}=\qty{6,7}{g}$, soit $\eta\approx88~\%$ —
mais le solide resterait impur. Pour un médicament, la \res{\text{pureté prime sur le rendement}} : on accepte la
perte.
\end{questions}
\end{corrige}
\end{blocLivre}
\begin{blocLivre}{cor-12-15}
\begin{corrige}{15}{Une estérification plus verte}
\demarche{L'économie d'atomes ne dépend que de l'équation ; le rendement, lui, dépend du procédé.}
\begin{questions}
\item $\ce{CH3COOH + C2H5OH <=> CH3COOC2H5 + H2O}$ :
$\mathrm{EA}=\dfrac{88{,}1}{60{,}1+46{,}1}=\res{83~\%}$. Elle ne dépend que de l'équation, donc
\res{\text{pas du procédé}}.
\item Le tamis moléculaire retire l'eau : $Q_\mathrm{r}$ reste inférieur à $K$ et la transformation se poursuit vers
un état quasi total. La lipase, elle, ne fait qu'accélérer : \res{\text{c'est le retrait de l'eau}} qui élève le
rendement.
\item B travaille à 40~°C au lieu de 80~°C (moins d'énergie) et sans acide sulfurique concentré ni vapeurs portées à
ébullition : \res{\text{plus sobre et plus sûr}}.
\item Catalyseur \res{\text{réutilisable}} d'origine biologique, température modérée, absence de solvant — trois
principes de la chimie verte.
\end{questions}
\end{corrige}
\end{blocLivre}
